\section{Distinct positive-box comparisons}\label{app:positive-box-predecessors}
Theorem~\ref{thm:positive-box-upper} supersedes the upper bounds in this
appendix; the proofs are included because they use different estimates.
Throughout the appendix $\rho>1$, $t=\rho-1$, $\alpha=1/\rho$, and
$\eta=1-\alpha$. Work first on $[1,\rho]^n$. The law $I$ is independent
endpoint rounding and $O$ is independently fair endpoint orientation.
For one product write $C,V,P,O$ for its upper, lower, independent,
and orientation values, and $D_I=C-P,D_O=C-O$.

\subsection{The symmetric low/high argument}
\begin{proposition}[A coarse finite constant]\label{prop:coarse-box}
For positive monomials on a box whose aspect ratios are at most $\rho$,
\[
 \tgap\le K(\rho)\hgap,\qquad
 K(\rho)=2(1+1/\rho)+\max\{4(1+\rho),8\eta^{-3}\}.
\]
In particular, with $R=\max\{4,\max_i u_i/\ell_i\}$,
$\tgap\le(4R+6+2/R)\hgap$, which is $45/2$ for aspect ratios at most
four. Direct use of $\rho=2$ gives the weaker constant $K(2)=67$.
\end{proposition}
\begin{proof}
Classify a normalized mean as low when it is at most $1/2$. Put $q_i=p_i$
for lows and $q_i=1-p_i$ for highs, so $q_i\le1/2$. Divide the product's
values by $\rho^h$, where $h$ is the number of high coordinates.
Let $L(s),H(s)$ count the respective deviations at least $s$, and let
$Q_L,Q_H$ be their sums. Integrals in this paragraph are over $[0,1/2]$.
Define
\[
 \begin{array}{ll}
 C_L=1+\int(\rho^{L(s)}-1)\,ds,&P_L=\prod_{\rm low}(1+tq_i),\\
 C_H=1+\int(\alpha^{H(s)}-1)\,ds,&P_H=\prod_{\rm high}(1-\eta q_i).
 \end{array}
\]
Common-threshold low successes and high failures occupy opposite ends
of the interval. Thus $C=C_L+C_H-1$, $P=P_LP_H$, and
\begin{equation}\label{eq:coarse-components}
 D_I=(C_L-P_L)+(C_H-P_H)+J_I,
 \quad J_I=(P_L-1)(1-P_H)\ge0.
\end{equation}
All three terms are nonnegative. Set
$A=C_L-1-tQ_L$, $B=C_H-1+\eta Q_H$, and let $q_L,q_H$ be the respective
largest deviations, zero for empty groups. Integer-count inequalities give
\begin{equation}\label{eq:dispersion-coarse}
 A\ge t^2(Q_L-q_L),\qquad B\ge\eta^2(Q_H-q_H).
\end{equation}
For the first, use $(1+t)^l-1-tl\ge t^2(l-1)_+$; for the second,
$\alpha^l-1+\eta l$ has increments $\eta(1-\alpha^l)\ge\eta^2$
for $l\ge1$. Integrating proves both.
Conditional orientation expectations at either end of the interval are
$b_L=(1+t/2)^{L(s)}$ and $b_H=(1-\eta/2)^{H(s)}$, so
\[
 O=1+2\int(b_Lb_H-1),\quad
 D_O=D_{OL}+D_{OH}+J_O,\quad J_O=2\int(b_L-1)(1-b_H).
\]
Here $D_{OL},D_{OH}\ge0$ are the separate-group deficiencies.
Expanding in powers of $t$ gives
$(1+t)^l-2(1+t/2)^l+1\ge[(1+t)^l-1-tl]/2$, since each order-$j\ge2$
coefficient is multiplied by $1-2^{1-j}\ge1/2$. Therefore
\begin{equation}\label{eq:orientation-coarse}
 D_O\ge A/2+J_O,\qquad J_O\ge(t\eta/2)\min(q_L,q_H).
\end{equation}
The high separate-group deficiency $D_{OH}$ is nonnegative by
convexity of the integer power at the midpoint of $\alpha$ and one.

Write $D_{IH}=C_H-P_H$. If $Q_H\le1$, put $q=q_H$ and $R'=Q_H-q$.
Bonferroni's bound gives
\[
 D_{IH}\ge B-\eta^2\sum_{i<j}q_iq_j
 \ge B-\eta^2(qR'+R'^2/2)
 \ge(1-q-R'/2)B\ge B/4.
\]
Here $B\ge\eta^2R'$, $q\le1/2$, and $q+R'\le1$; no division by
$R'$ is needed. If $Q_H\ge1$, order deviations decreasingly.
At a fixed sum, the linear common-threshold high value is minimized
by filling the largest-weight coordinates to their cap $1/2$ and at
most one remainder. Convex interpolation then gives
\[
 C_H\ge\tfrac12+\tfrac12\alpha^{2Q_H},\qquad P_H\le e^{-\eta Q_H}.
\]
The difference $g(Q)=1/2+\alpha^{2Q}/2-e^{-\eta Q}$ is nondecreasing
for $Q\ge1$. Indeed $-\log\alpha\le\eta/\alpha$ and
$\alpha^{2Q-1}\le e^{-\eta Q}$ imply
$g'(Q)\ge\eta[e^{-\eta Q}-\alpha^{2Q-1}]\ge0$.
The fourth-order exponential remainder yields
\[
 D_{IH}\ge g(1)=1-\eta+\eta^2/2-e^{-\eta}
 \ge\eta^3/6-\eta^4/24\ge\eta^3/8.
\]

Let $\phi_\rho$ interpolate $\rho^k$ between every consecutive integer,
including negative integers. The exact normalized lower envelope is
$V=\phi_\rho(Q_L-Q_H)$, because the original count sum is shifted by
integer $h$. Supporting slopes at zero are $\eta$ and $t$, so
\begin{equation}\label{eq:coarse-target}
 C-V\le A+B+t\eta\min(Q_L,Q_H).
\end{equation}
For $Q_H\le1$, use
$\min(Q_L,Q_H)\le\min(q_L,q_H)+A/t^2+B/\eta^2$ to obtain
\[
 T\le(1+1/\rho)A+(1+\rho)B+2J_O
 \le2(1+1/\rho)D_O+4(1+\rho)D_I.
\]
For $Q_H\ge1$, $T\le C_L=1+A+tQ_L$ and
$J_I\ge tQ_L(1-e^{-\eta})\ge t\eta Q_L/2$.
The preceding bounds therefore imply
\[
 T\le2D_O+8\eta^{-3}D_{IH}+2\eta^{-1}J_I
      \le2D_O+8\eta^{-3}D_I.
\]
Combining the cases proves the per-product inequality with the two
coefficients in $K(\rho)$. Mixing the global laws proportionally, undoing
the product's normalization, and summing proves the common-aspect bound.
The positive affine substitution of
Theorem~\ref{thm:positive-box-upper} transfers it to unequal aspects.
For $R\ge4$, $8(1-1/R)^{-3}\le512/27<20\le4(1+R)$, giving
$K(R)=4R+6+2/R$. For $\rho=2$ the two coefficients are $3$ and $64$.
\end{proof}
The high-mass argument in fact allows $8\eta^{-3}$ to be replaced by
$1/g(1)$. Indeed $g(1)\le1-e^{-\eta}$, since their difference is
$\eta-\eta^2/2\ge0$, so the same coefficient controls both the constant
and the $J_I$ term in $C_L$. This gives the sharper finite constant
$2(1+1/\rho)+\max\{4(1+\rho),1/g(1)\}$.

The tempting all-high simplification $T\le2D_I$ is false. On $[1,2]^3$
take failure probabilities $(1/2,49/100,1/100)$ and divide values by
$2^3$. The exact values are
\[
 C=501/800,\quad V=1/2,\quad P=90147/160000,
 \quad T=20200/160000>2D_I=20106/160000.
\]
The boundary mean $1/2$ is allowed in this all-high calculation. The
example does not contradict the mixed-law inequality.

\subsection{The asymmetric split}
\begin{proposition}[An asymmetric leading-order bound]\label{prop:asymmetric-box}
For $\rho\ge64$, positive products on boxes of aspect ratios at most
$\rho$ satisfy
\[
 \tgap\le\left(2+\frac{\rho+1}{1-3/\sqrt\rho}\right)\hgap.
\]
\end{proposition}
\begin{proof}
Put $\delta=\rho^{-1/2}\le1/8$ and $b=(\rho+1)/(1-3\delta)$.
Now classify a normalized mean as low if $p_i\le1-\delta$ and high
otherwise, so high deviations are $q_i=1-p_i<\delta$. Normalize by
$\rho^h$ again, and redefine the counts as
$L(s)=\#\{i\text{ low}:p_i\ge s\}$ and
$H(s)=\#\{i\text{ high}:q_i\ge s\}$, with
$Q_L=\sum_{\rm low}p_i$ and $Q_H=\sum_{\rm high}q_i$.
The separate groups' common-threshold upper values and independent
values are now
\[
 \begin{array}{ll}
 C_L=1+\int_0^{1-\delta}(\rho^{L(s)}-1)\,ds,
 &P_L=\prod_{\rm low}(1+tp_i),\\
 C_H=1+\int_0^\delta(\alpha^{H(s)}-1)\,ds,
 &P_H=\prod_{\rm high}(1-\eta q_i).
 \end{array}
\]
Low successes and high failures under common-threshold rounding are
disjoint, so \eqref{eq:coarse-components}, the definitions
$A=C_L-1-tQ_L$, $B=C_H-1+\eta Q_H$, and
\eqref{eq:coarse-target} remain valid. The dispersion bounds
\eqref{eq:dispersion-coarse} hold with maxima $a_L,a_H$ of low means
and high failures; they did not require low means at most $1/2$.
A different estimate now replaces $D_{OL}\ge A/2$:
\[
 D_{IL}=C_L-P_L\ge\delta A.
\]
To prove it, expand the low physical product. Constant and affine terms
cancel. Every subset of degree at least two has independent product at
most $(1-\delta)$ times its minimum mean, and the common upper values of
these subsets sum to exactly $A$.

For $s\in[0,\delta]$ and $U=s$ or $U=1-s$, conditional low expectations
are $1+(t/2)\ind\{s\le p_i\}$, while high normalized expectations
are $1-(\eta/2)\ind\{s\le q_i\}$. Outside these two end strips
the conditional high product is one. Conditional independence of the
orientation coins gives $D_O=D_{OL}+D_{OH}+J_O\ge J_O$, where
\[
 J_O=2\int_0^\delta
 \left(\prod_{\rm low}[1+(t/2)\ind\{s\le p_i\}]-1\right)
 \left(1-\prod_{\rm high}[1-(\eta/2)\ind\{s\le q_i\}]\right)ds
 \ge(t\eta/2)\min(a_L,a_H).
\]
Combining this with dispersion and the target inequality yields
$T\le(1+1/\rho)A+(1+\rho)B+2J_O$.
If $Q_H\le4\delta$, the same Bonferroni argument with $q=a_H\le\delta$
and $R'=Q_H-q\le4\delta$ gives
$D_{IH}\ge(1-q-R'/2)B\ge(1-3\delta)B$.
Since $(1+1/\rho)/\delta\le b$, it follows that
$T\le bD_I+2D_O$.

If $Q_H\ge4\delta$, the common high failure set has measure at most
$\delta$, so $C_H\ge1-\delta$. Also
$P_H\le e^{-4\eta\delta}\le1-2\delta$, by
$e^{-z}\le1-z+z^2/2$ and $4\eta-8\eta^2\delta\ge3-1=2$.
Hence $D_{IH}\ge\delta$ and $J_I\ge\delta tQ_L$, giving
$D_I\ge\delta(A+1+tQ_L)=\delta C_L$.
As $T\le C_L$, again $T\le bD_I+2D_O$.
Mix the two global laws with weights proportional to $b,2$, sum over
products, and use positive affine transfer for unequal aspects.
\end{proof}
The constant in Proposition~\ref{prop:asymmetric-box} is
$\rho+3\sqrt\rho+O(1)$; Theorem~\ref{thm:positive-box-upper} improves
the additive error to two.

\subsection{Optimality within a fixed two-law mixture}
\begin{proposition}[A limitation of fixed independent/fair-orientation mixtures]
\label{prop:fixed-mixture-optimal}
Fix $\rho>1$ and an orientation weight $w\in[0,1]$. For mixtures of
$I$ and independently fair $O$, the best uniform per-physical-monomial
captured fraction over all dimensions is at most
$\min\{w/2,(1-w)/\rho\}$. Optimizing the weight gives exactly
$1/(\rho+2)$, attained at $w=2/(\rho+2)$.
\end{proposition}
\begin{proof}
For a bilinear product at normalized means $(\epsilon,1-\epsilon)$,
$0<\epsilon<1/2$, deleting affine parts gives
$D_I/T=\epsilon$ and $D_O/T=1/2$. Letting $\epsilon\downarrow0$
bounds the mixture fraction by $w/2$.
For the other obstruction use means $\epsilon$ and $k$ copies of
$1-\epsilon/k$, where $k\ge1$. Divide physical values by $\rho^{k+1}$,
and write $\alpha=1/\rho$, $\eta=1-\alpha$, $\gamma=(1+\alpha)/2$.
The sum of success means is $k$, giving $V=\alpha$.
Common-threshold and fair-orientation integration, split at
$\epsilon/k$ and $\epsilon$, and direct independent multiplication give
\[
 \begin{split}
 C&=\alpha+\epsilon[\eta-\alpha(1-\alpha^k)/k],\\
 P&=\alpha[1+(\rho-1)\epsilon][1-\eta\epsilon/k]^k,\\
 O&=\alpha+\epsilon[\eta-2\gamma(1-\gamma^k)/k].
 \end{split}
\]
In the orientation law, the two end intervals of length
$\epsilon/k$ have normalized expected product $\gamma^{k+1}$,
the two next intervals of total length $2(\epsilon-\epsilon/k)$
have value $\gamma$, and the middle interval of length $1-2\epsilon$
has value $\alpha$; these yield the stated $O$.
For each fixed $k$, divide deficiencies by $\epsilon$ and let
$\epsilon\downarrow0$. Then
\[
 \begin{split}
 T/\epsilon&=\eta-\alpha(1-\alpha^k)/k,\\
 D_I/\epsilon&\longrightarrow\alpha\eta-\alpha(1-\alpha^k)/k,\\
 D_O/\epsilon&=[2\gamma(1-\gamma^k)-\alpha(1-\alpha^k)]/k.
 \end{split}
\]
The denominator coefficient is at least $\eta^2>0$ because
$1-\alpha^k\le k\eta$. Now let $k\to\infty$ to obtain
$D_I/T\to1/\rho$ and $D_O/T\to0$, bounding the mixture by
$(1-w)/\rho$. The maximum of the minimum of these two bounds is
$1/(\rho+2)$ at their intersection. Lemma~\ref{lem:box-coefficient}
with fair orientations supplies the matching guarantee.
\end{proof}
The proposition concerns only the per-monomial guarantee of a fixed
two-law mixture over unrestricted dimensions. It does not cover other
laws, coefficient-dependent mixtures, or the optimal joint law for a
sum of monomials. In particular, it does not prove
$C_{\rm box}(\rho)=\rho+2$ and does not conflict with the
finite-dimensional balanced-orientation refinement.
